\documentclass{book}
\usepackage{a4}
\usepackage[czech]{babel}
\usepackage[utf8]{inputenc}
\usepackage[T1]{fontenc}
\usepackage{textcomp}
\usepackage{subcaption}
\usepackage{gensymb}
\usepackage{pdfpages}
\usepackage{float}
\usepackage{subfiles}
%podminene formatovani
\usepackage{ifthen}
\usepackage{array}

\usepackage{tikz}
%\usepackage{background}
%\backgroundsetup{contents=Pracovní verze}

\usetikzlibrary{calligraphy}

%\backgroundsetup{
%    scale=2.5,
%    angle=90,
%    firstpage = true, 
%    opacity=0.1, 
%    contents = {
%        \begin{tikzpicture}[opacity=.2]
%    			\clip (1,-15) rectangle (11,15);
%    			\foreach \x in {0,0.05,...,2}{
%      			\draw (0,\x) .. controls (4-\x,-4-3*\x) and (11-2*\x,4-3*\x) .. (12,-5+3*\x);
%      			\draw (0,\x-1) .. controls (4-\x,\x-2) and (11-2*\x,-4-3*\x) .. (12,-3+3*\x);
%      			\draw (0,\x-3) .. controls (4-\x,\x-2) and (11-4*\x,-1-4*\x) .. (12,-5+3.5*\x);
%    			}
%  		\end{tikzpicture}
%    }
%} 

%udaje o zhotoviteli
\newcommand{\zhotovitel}{secureWORK s.r.o.}
\newcommand{\zhotovitelIC}{06694497}
\newcommand{\zhotovitelUlice}{Bělehradská 858/23}
\newcommand{\zhotovitelMesto}{Praha}
\newcommand{\zhotovitelPSC}{120 00}

%udaje o zadateli
\newcommand{\objednatel}{}
\newcommand{\objednatelEmail}{}

\newcommand{\objednatelInstalaceUlice}{}
\newcommand{\objednatelInstalaceMesto}{}
\newcommand{\objednatelInstalacePSC}{}


%udaje o elektrarne
\newcommand{\vykonFVE}{}
\newcommand{\kapacitaBaterii}{}
\newcommand{\stridac}{}
\newcommand{\baterie}{}
\newcommand{\wallbox}{}
\newcommand{\pocetPanelu}{} %umisti mezeru za cislo!
\newcommand{\pocetVstupuPoleA}{} %pokud je pouzit jen 1 vstup = 1, pokud oba = 2
\newcommand{\pocetPaneluPoleAA}{}
\newcommand{\pocetPaneluPoleAB}{}
\newcommand{\pocetPaneluPoleB}{} %umisti mezeru za cislo!

\newcommand{\zjisteniceMenic}{}
\newcommand{\zjisteniceWallbox}{}
\newcommand{\zjistenicePanely}{}

\newcommand{\zjisteniceDCBOX}{}

\newcommand{\zjisteniceRevize}{}

\newcommand{\zjisteniceTermokontrola}{Termografická kontrola komponent \newline Střídač: \newline Baterie: \newline DCBOX: \newline Rozvaděč: \newline Kabelové vedení:}

\newcommand{\zjisteniceACDC}{}

\newcommand{\zjisteniceChyby}{}

\newcommand{\zjisteniceOstatni}{}


\newcommand{\mistoMereni}{}
\newcommand{\datumMereni}{}


\begin{document}
 
	\chapter*{Periodická kontrola FVE1}
	

	\subsection*{Provádějící}		
	\begin{tabbing}
	\hspace{2em} \= \textbf{Firma:}\hspace{6em}  \= \textbf{\zhotovitel}, \zhotovitelIC / CZ\zhotovitelIC \\[0.3cm]
	\> \textbf{Adresa:} \>  \zhotovitelUlice\\[0.3cm]
	\> \textbf{PSČ:} \>  \zhotovitelPSC\\[0.3cm]
	\> \textbf{E-mail:} \>   kobrle@securework.cz\\
	\end{tabbing}
	

	\subsection*{Místo instalace}	
	
	\begin{tabbing}
	\hspace{2em} \= \textbf{Adresa:} \hspace{10em}  \= \objednatelInstalaceUlice \objednatelInstalaceMesto \\[0.3cm]
	\> \textbf{PSČ:} \>  \objednatelInstalacePSC \\[0.3cm]
	\> \textbf{Majitel:} \>   \objednatel \\[0.3cm]
	\> \textbf{E-mail:} \>   \objednatelEmail \\[0.3cm]
	\end{tabbing}
	
	\subsection*{Konfigurace instalace}	
	
	\begin{tabbing}
	\hspace{2em} \= \textbf{Výkon panelů:} \hspace{6em}  \= \vykonFVE \\[0.3cm]
	\> \textbf{Konfigurace stringů:} \>  \pocetPaneluPoleAA  \pocetPaneluPoleB \\[0.3cm]
	\> \textbf{Střídač:} \>  \stridac \\[0.3cm]
	\> \textbf{Akumulátory:} \>   \baterie \,  \kapacitaBaterii \\[0.3cm]
	\> \textbf{Wallbox:} \>   \wallbox \\[0.3cm]
	\end{tabbing}
	
	\newpage
	
	\section*{Kontrola}

	Obsah kontroly byl proveden v rozsahu minimálně dle rozsahu v Tab. \ref{tabulka_kontrolni_seznam} a výsledky zaznamenány v Tab. \ref{tabulka_protokol_kontrolniho_seznamu}. 
	
	Případné další zaznamenané nedostatky fotovoltaické elektrárny:
	
	\begin{center}
	\vskip.2in
	............................................................................................................
	\vskip.2in
	............................................................................................................
	\vskip.2in
	............................................................................................................
	\vskip.2in
	............................................................................................................
	\vskip.2in
	............................................................................................................
	\vskip.2in
	............................................................................................................
	\vskip.2in
	\end{center}
	
	
	Podpisem prováděcí technik potvrzuje provedené úkony ke dni kontroly.
	\vskip0.5in
	\begin{tabbing}
	
	\hspace{18em} \= V \mistoMereni \, dne: \datumMereni \\[2.5cm]
	
	\hspace{21em} \= ............................................. \\[0.5cm]
	
	\hspace{25em} \= \textbf{TECHNIK} \\[1.5cm]
	\end{tabbing}
	

	\begin{table}
	\begin{tabular}{|c|p{.3\textwidth}|p{.5\textwidth}|}
	\hline 
	\textbf{PERIODA} & \textbf{KOMPONENTA} & \textbf{ČINNOST} \\ 
	\hline 
	DENNĚ & Měnič & Je v provozu bez indikace poruchy? \\ 
	\hline 
	MĚSÍČNĚ & Kontrola výnosu & Pravidelně zaznamenávejte stavy elektroměru! \\ 
	\hline 
	PŮL ROČNĚ & Plocha panelů & 
	\begin{itemize}
	\item V případě silného znečištění listím, ptačím trusem , znečištěným ovzduším nebo něčím jiným - proveďte čištění.
	\item Jsou moduly správně upevněny?
	\item Je plocha generátoru pod mechanickým pnutím?
	\end{itemize} \\ 
	\hline 
	ROČNĚ & Termografická kontrola & 
	\begin{itemize}
	\item Termální foto panelů.
	\item Termální foto DCBOXU + vodičů.
	\item Termální foto měniče + aku.
	\item Termální rozvaděčů a svorkovnic.
	\end{itemize} \\ 
	\hline 
	ROČNĚ & Rozváděč generátoru / DCBOXu & 
	\begin{itemize}
	\item Vniklý hmyz / vlhkost?
	\item Zkontrolujte pojistky.
	\item Utažení svorek vodičů.
	\end{itemize} \\ 
	\hline 
	ROČNĚ & Přepěťová ochrana  & 
	\textbf{Zkontrolujte i po bouřkách!}
	\begin{itemize}
	\item Je přepěťová ochrana v pořádku?
	\end{itemize} \\ 
	\hline 
	ROČNĚ & Vedení  & 
	\begin{itemize}
	\item Seškvařená místa, přerušená izolace či jiná poškození?
	\item Zkontrolujte spojky.
	\end{itemize} \\ 
	\hline 
	1x za 3-4 ROKY & Opakování měření pro uvádění do provozu  & 
	Kontrolu provede odborník. \\ 
	\hline 
	1x za 3-4 ROKY & Měnič (venkovní)  & 
	Kontrola proti vniknutí vlhkosti, zvířat či jiných těles. Kontrolu provede odborník. \\ 
	\hline 
	PŘI PODEZŘENÍ & Moduly  & 
	Měření charakteristiky, termografická kontrola nebo analýza funkce provedená odborníkem. \\ 
	\hline 
	PŘI PODEZŘENÍ & Rozváděč generátoru  & 
	Zkontrolujte DC pojistky všech stringů. \\ 
	\hline 
	PŘI PODEZŘENÍ & Ochrany na AC straně  & 
	\begin{itemize}
	\item Kontrola jištění vedení.
	\item AC pojistky a proudový chránič.
	\end{itemize}
	 \\ 
	\hline 
	
	\end{tabular} 
	\caption{Kontrolní seznam pro údržbu a servis}
	\label{tabulka_kontrolni_seznam}
	\end{table}
	
	\begin{table}
	\begin{tabular}{| m{4cm} | p{.2\textwidth} | p{.5\textwidth} |}
	\hline 
	\textbf{KOMPONENTA} & \textbf{VÝSLEDEK} & \textbf{POZNÁMKA} \\ 
	\hline 
	  MĚNIČ & \begin{tabular}{@{}c@{}}V POŘÁDKU \\ S VÝHRADOU \\ ZÁVADA\end{tabular} & \zjisteniceMenic \\[0.5cm]
	\hline 
	  WALLBOX & \begin{tabular}{@{}c@{}}V POŘÁDKU \\ S VÝHRADOU \\ ZÁVADA\end{tabular} & \zjisteniceWallbox \\[0.5cm]
	\hline 
	  FV PANELY & \begin{tabular}{@{}c@{}}V POŘÁDKU \\ S VÝHRADOU \\ ZÁVADA\end{tabular} & \zjistenicePanely \\[0.5cm]
	\hline 
	  DCBOX & \begin{tabular}{@{}c@{}}V POŘÁDKU \\ S VÝHRADOU \\ ZÁVADA\end{tabular} & \zjisteniceDCBOX \\[0.5cm]
	\hline 
	  REVIZNÍ MĚŘENÍ & \begin{tabular}{@{}c@{}}V POŘÁDKU \\ S VÝHRADOU \\ ZÁVADA\end{tabular} & \zjisteniceRevize \\[0.5cm]
	\hline 
	  TERMOGRAFICKÁ KONTROLA & \begin{tabular}{@{}c@{}}V POŘÁDKU \\ S VÝHRADOU \\ ZÁVADA\end{tabular} & \zjisteniceTermokontrola \\[0.5cm]
	\hline 
	  ROZVADĚČE AC / DC & \begin{tabular}{@{}c@{}}V POŘÁDKU \\ S VÝHRADOU \\ ZÁVADA\end{tabular} & \zjisteniceACDC \\[0.5cm]
	\hline 
	  KONTROLA CHYB & \begin{tabular}{@{}c@{}}V POŘÁDKU \\ S VÝHRADOU \\ ZÁVADA\end{tabular} & \zjisteniceChyby \\[3.5cm]
	\hline 
	  OSTATNÍ & \begin{tabular}{@{}c@{}}V POŘÁDKU \\ S VÝHRADOU \\ ZÁVADA\end{tabular} & \zjisteniceOstatni \\[2.5cm]
	\hline 
	
	\end{tabular} 
		
	\caption{Protokol údržby a servisu FVE}
	\label{tabulka_protokol_kontrolniho_seznamu}
	\end{table}
	
	%\newpage
	%\includepdf[pages=-, angle=0]{revize.pdf}
	
\end{document}